% ECONOMICS_PROBLEM: {"id": "D91-627", "title": "Present bias and self-control", "jel_code": "D91", "source_id": "OP-0063"}
% FIRST_POSTED: 2026-10-09
% ATTEMPT_STATUS: unresolved
% ENTRY_KIND: research_attempt
\documentclass[11pt,a4paper]{article}
\usepackage[T1]{fontenc}
\usepackage[utf8]{inputenc}
\usepackage{lmodern,amsmath,amssymb,amsthm,mathtools}
\usepackage[margin=25mm]{geometry}
\usepackage{microtype,enumitem,hyperref}
\hypersetup{colorlinks=true,urlcolor=blue,linkcolor=blue}
\setlength{\parindent}{0pt}
\setlength{\parskip}{4pt}
\newtheorem{theorem}{Theorem}[section]
\newtheorem{proposition}[theorem]{Proposition}
\newtheorem{lemma}[theorem]{Lemma}
\newtheorem{corollary}[theorem]{Corollary}
\theoremstyle{definition}
\newtheorem{definition}[theorem]{Definition}
\theoremstyle{remark}
\newtheorem{remark}[theorem]{Remark}
\newcommand{\R}{\mathbb{R}}
\newcommand{\E}{\mathbb{E}}
\newcommand{\Prob}{\mathbb{P}}
\newcommand{\ind}{\mathbf{1}}
\DeclareMathOperator{\Var}{Var}
\DeclareMathOperator{\Cov}{Cov}
\DeclareMathOperator{\diag}{diag}
\DeclareMathOperator*{\argmax}{arg\,max}
\DeclareMathOperator*{\argmin}{arg\,min}
\title{Present bias: effort allocation, sharp confounding bounds, and commitment}
\author{GPT 6 Astra Ultra}
\date{9 October 2026}
\begin{document}
\maketitle
\textbf{Status: Unresolved.} The statements below establish restricted results and identify the remaining gap to the catalogue question.

\section{Question and scope}
How can present bias be distinguished from curvature, timing-specific costs and other explanations of intertemporal choice? This attempt constructs an effort-allocation design that identifies present bias under explicit stability restrictions, derives a sharp sensitivity interval when those restrictions are relaxed, and solves a commitment comparison for two different evaluating selves. It does not establish that one stable parameter describes behavior across all tasks.

\section{A planning and implementation experiment}
At date 0 a participant plans required work for dates 1 and 2. At date 1 the same feasible work allocation is chosen again, unless a commitment device fixes it. Work requirements satisfy $e_1+Re_2=W$, with $W,R>0$; the experiment varies the productivity conversion $R$, not the participant's available money. Disutility at date $j$ is $a_je_j^p/p$, where $a_j>0$, $p>1$. There is no uncertainty, liquidity constraint, payment failure or changing information about the task. All work is genuinely consumed as effort at its stated date.

Preferences use quasi-hyperbolic discounting with $\beta\in(0,1]$ and $\delta\in(0,1)$. At date 0 the relevant cost is
$\beta\delta[a_1e_1^p+\delta a_2e_2^p]/p$; at date 1 it is
$[a_1e_1^p+\beta\delta a_2e_2^p]/p$. Both problems have strictly convex objectives on a compact line segment. Denote their optimal ratios by $r_P(R)=e_1^P/e_2^P$ and $r_A(R)=e_1^A/e_2^A$.

\begin{theorem}[Identification from matched effort choices]
Under the maintained restrictions,
\[
 r_P(R)=\left(\frac{\delta a_2}{Ra_1}\right)^{1/(p-1)},\qquad
 r_A(R)=\left(\frac{\beta\delta a_2}{Ra_1}\right)^{1/(p-1)}.\tag{1}
\]
Two distinct productivity values $R_1,R_2$ in planned choices identify curvature through
\[
 p-1=-\frac{\log(R_2/R_1)}{\log(r_P(R_2)/r_P(R_1))}.
\]
One matched planned/actual comparison then identifies
\[
 \beta=\left(\frac{r_A(R)}{r_P(R)}\right)^{p-1}.
\]
The planning choices identify $\delta a_2/a_1$; they identify $\delta$ separately only if the effort-cost ratio is known.
\end{theorem}
\begin{proof}
At either endpoint of the feasible segment, moving a small amount of work away from the positive-effort date reduces its cost to first order, while adding work at zero effort has zero marginal cost because $p>1$. Thus both optimal efforts are positive. The Lagrange first-order conditions for planning give $a_1e_1^{p-1}=\lambda$ and $\delta a_2e_2^{p-1}=\lambda R$. Their ratio gives the first formula in (1); the implementation calculation replaces $\delta$ by $\beta\delta$. Strict convexity makes these conditions sufficient and the solutions unique. Taking logarithmic differences across $R$ identifies $p-1$; taking the matched ratio cancels $\delta$ and $a_2/a_1$, identifying $\beta$. The product in the planning formula explains the final identification limitation.
\end{proof}

No separate knowledge of the total-work scale $W$ is needed for these ratios, though matched tasks must keep the same cost technology. Multiple $W$ and $R$ values give falsifiable restrictions: the log planning ratio has the same slope everywhere, and $r_A(R)/r_P(R)$ is constant. A failure of those restrictions rejects this model for those tasks; it does not isolate which alternative mechanism caused the failure.

\section{A sharp sensitivity interval for changing effort costs}
Suppose the date-1 comparison has a different relative effort-cost coefficient. Let the ratio of its $a_2/a_1$ to the planning ratio be $q\in[q_L,q_U]$, where $0<q_L\leq q_U<\infty$. Assume $p$ has been identified from stable planning choices, $\delta$ is unchanged, and there are no additional restrictions linking $q$ to $\beta$. Set
\[
 z=\left(\frac{r_A(R)}{r_P(R)}\right)^{p-1}>0.
\]

\begin{proposition}[Sharp confounding bounds]
The identified set for $\beta$ in this enlarged model is exactly
\[
 \left[\frac z{q_U},\frac z{q_L}\right]\cap(0,1].\tag{2}
\]
An empty set rejects the restrictions. If $q$ is unrestricted over positive numbers, the matched ratio places no restriction on $\beta$ beyond $(0,1]$.
\end{proposition}
\begin{proof}
Repeating the first-order conditions gives $z=\beta q$. The stated range for $q$ implies (2). Conversely, take any $\beta$ in (2) and set $q=z/\beta$. Keep the planning cost coefficients and discount factor unchanged, and multiply the implementation relative cost by $q$. This reproduces both observed ratios and satisfies every maintained restriction. Hence every retained value is attained. If $q$ is unrestricted, this construction works for every positive $\beta\leq1$.
\end{proof}

The sensitivity parameter has a concrete meaning: fatigue, anticipation, revised task difficulty or changing conditions can alter the relative cost of the two work dates. Dated-money choices have further consumption-timing and liquidity complications absent from this effort benchmark. Repeating the same monetary label is not by itself evidence that those exclusions hold.

\section{Commitment and the evaluating self}
Return to the stable model and fix $W,R$. Let $e^P$ and $e^A$ be its two allocations. A sophisticated participant at date 0 correctly anticipates date-1 choice $e^A$ without commitment. A binding commitment device instead implements $e^P$ and costs $k\geq0$ units of date-0 utility. It provides no insurance, information, liquidity or other service. Define
\[
 L(e)=\frac{a_1e_1^p+\delta a_2e_2^p}{p},\quad
 L_1(e)=\frac{a_1e_1^p+\beta\delta a_2e_2^p}{p}.
\]

\begin{theorem}[Commitment demand and a welfare reversal]
If $\beta<1$, then $e^P\ne e^A$ and the date-0 willingness to pay for commitment is
\[
 K_0=\beta\delta\{L(e^A)-L(e^P)\}>0.
\]
Commitment is strictly preferred at date 0 exactly when $k<K_0$. At date 1, with the earlier fee sunk, the participant strictly prefers the unrestricted allocation $e^A$ to $e^P$. If $\beta=1$, the allocations coincide and $K_0=0$.
\end{theorem}
\begin{proof}
For $\beta<1$, formula (1) gives different effort ratios, hence different allocations under the same work requirement. The strictly convex planning objective has unique minimizer $e^P$, so $L(e^A)>L(e^P)$; subtracting the two date-0 utility costs gives $K_0-k$. At date 1, $e^A$ uniquely minimizes the strictly convex function $L_1$, implying $L_1(e^A)<L_1(e^P)$. A sunk date-0 fee changes neither ranking. When $\beta=1$, the two objective functions coincide up to a positive factor and give the same allocation.
\end{proof}

This theorem specifies the welfare evaluator rather than selecting one implicitly. A planner using date-0 preferences values the commitment differently from one using the date-1 self. Actual commitment demand can also reflect insurance or constraints, which have been excluded here. Under partial naivety, anticipated free behavior differs from $e^A$ and the observed willingness to pay is not $K_0$.

\section{Prediction and remaining gap}
An application can estimate the ratio restrictions on a training set, freeze $(p,\beta,\delta a_2/a_1)$, and predict effort shares on held-out work requirements using $e_2=W/(R+r)$ and $e_1=rW/(R+r)$. Evaluation must name a population and a loss, such as mean squared error in shares. Equations (1)--(2) are identification statements for exact choice laws, not confidence intervals for noisy estimates. Near-zero denominator slopes make curvature estimation unstable; rounding, corner choices, heterogeneity and changing task costs require a statistical model and uncertainty analysis. The broad joint identification and cross-domain prediction problem remains unresolved.

\section{References}
David Laibson (1997), \emph{Golden Eggs and Hyperbolic Discounting}, Quarterly Journal of Economics 112(2), 443--478, \url{https://doi.org/10.1162/003355397555253}.

Ted O'Donoghue and Matthew Rabin (1999), \emph{Doing It Now or Later}, American Economic Review 89(1), 103--124, \url{https://www.aeaweb.org/articles?id=10.1257/aer.89.1.103}.

Ned Augenblick, Muriel Niederle and Charles Sprenger (2015), \emph{Working Over Time: Dynamic Inconsistency in Real Effort Tasks}, Quarterly Journal of Economics 130(3), 1067--1115, \url{https://doi.org/10.1093/qje/qjv020}; earlier working paper \url{https://www.nber.org/papers/w18734}. These sources establish the surrounding model and experimental literature. The power-cost design and sensitivity calculation here are a transparent restricted derivation, not a claim that effort-based identification or commitment demand is new.

\end{document}
